% P-002: The Arithmetic Balance Ratio R(n) = sigma(n)phi(n)/(n*tau(n))
%         and Characterizations of the Number 6
%
% Target: American Mathematical Monthly (or similar expository journal)
% Status: Draft v1.0  (2026-03-26)

\documentclass[12pt]{article}

\usepackage{amsmath}
\usepackage{amssymb}
\usepackage{amsthm}
\usepackage{mathtools}
\usepackage{array}
\usepackage{booktabs}
\usepackage{hyperref}
\usepackage{geometry}
\geometry{margin=1in}

% Theorem environments
\newtheorem{theorem}{Theorem}[section]
\newtheorem{corollary}[theorem]{Corollary}
\newtheorem{lemma}[theorem]{Lemma}
\newtheorem{proposition}[theorem]{Proposition}
\newtheorem{definition}[theorem]{Definition}
\newtheorem{remark}[theorem]{Remark}

% Shorthand macros
\newcommand{\sig}{\sigma}
\newcommand{\ph}{\varphi}
\newcommand{\SpecR}{\mathrm{Spec}_R}
\newcommand{\gcd}{\mathrm{gcd}}
\newcommand{\lcm}{\mathrm{lcm}}
\newcommand{\Z}{\mathbb{Z}}
\newcommand{\N}{\mathbb{N}}
\newcommand{\Q}{\mathbb{Q}}
\newcommand{\R}{\mathbb{R}}
\newcommand{\C}{\mathbb{C}}

\title{The Arithmetic Balance Ratio $R(n) = \dfrac{\sigma(n)\,\varphi(n)}{n\,\tau(n)}$\\[6pt]
       and Characterizations of the Number $6$}

\author{%
  % Author information omitted for review
}

\date{March 2026}

\begin{document}

\maketitle

\begin{abstract}
We introduce the \emph{arithmetic balance ratio}
\[
  R(n) \;=\; \frac{\sigma(n)\,\varphi(n)}{n\,\tau(n)},
\]
where $\sigma$, $\varphi$, and $\tau$ denote the sum-of-divisors, Euler totient,
and number-of-divisors functions, respectively.  We develop a complete theory
of $R$ and use it to give several new characterizations of the number $6$.

Our main results are as follows.  (A) $R$ is multiplicative, with
$R(p^a) = (p^{a+1}-1)/(p(a+1))$ on prime powers.  (B) The \emph{kernel}
$\ker(R) := \{n\in\N : R(n)=1\}$ equals $\{1,6\}$, and $6$ acts as an
identity element for $R$ under coprime multiplication: $R(6n)=R(n)$ whenever
$\gcd(n,6)=1$.  (C) The pair $(2,3)$ is the unique pair of primes $p\le q$
such that $R(p)R(q)=1$.  (D) The \emph{$R$-spectrum}
$\SpecR := \{R(n):n\ge 1\}$ is discrete and satisfies the
\emph{Gap Theorem}: $\SpecR = \{3/4\}\cup\{1\}\cup[7/6,+\infty)$, so the
intervals $(3/4,1)$ and $(1,7/6)$ are free of $R$-values.  (E) The
\emph{Topological Master Formula} $\sigma(n)\,\varphi(n)\,f(n)=1$ holds if and
only if $n=6$, where $f(n)=\delta^+(R(n))\cdot\delta^-(R(n))$ is the product
of the one-sided spectral gaps at $R(n)$.  (F) The identity
$\varphi(n)/\tau(n)+\tau(n)/\sigma(n)+1/n=1$ holds if and only if $n=6$,
realizing $\tfrac12+\tfrac13+\tfrac16=1$ as an equation in arithmetic
functions.  We also compute the Dirichlet series $F(s)=\sum_{n=1}^\infty
R(n)/n^s = \zeta(s)\zeta(s+1)$ and show that $R(P_k)$ is an integer for every
even perfect number $P_k$.

The entire theory traces to one Diophantine fact: $(p^2-1)(q^2-1)=4pq$ has the
unique prime solution $(p,q)=(2,3)$.
\end{abstract}

\tableofcontents
\bigskip

% ------------------------------------------------------------------
\section{Introduction}
\label{sec:intro}
% ------------------------------------------------------------------

For a positive integer $n$, let $\sigma(n)$ denote the sum of its positive
divisors, $\varphi(n)$ the Euler totient (count of integers in $[1,n]$ coprime
to $n$), and $\tau(n)$ the number of positive divisors.  These three classical
arithmetic functions satisfy the well-known identities
\[
  \sigma(p^a)=\frac{p^{a+1}-1}{p-1},\quad
  \varphi(p^a)=p^{a-1}(p-1),\quad
  \tau(p^a)=a+1
\]
for prime powers, and all three are \emph{multiplicative}: $f(mn)=f(m)f(n)$
whenever $\gcd(m,n)=1$.

A number $n$ is \emph{perfect} if $\sigma(n)=2n$.  The identity
$\sigma(n)\varphi(n)=n\,\tau(n)$ characterizes $n\in\{1,6\}$ (see
OEIS A066068), and is the starting point for our work.

\begin{definition}
The \emph{arithmetic balance ratio} of a positive integer $n$ is
\[
  R(n) \;:=\; \frac{\sigma(n)\,\varphi(n)}{n\,\tau(n)}.
\]
\end{definition}

The ratio $R(n)$ measures how close $n$ is to satisfying $\sigma\varphi=n\tau$:
by definition, $R(n)=1$ if and only if $n\in\{1,6\}$.  Table~\ref{tab:Rvalues}
records $R(n)$ for $1\le n\le 20$, and already reveals the gap structure that
is the core of this paper.

\medskip
\noindent\textbf{Main Results.}
We state the main theorems here; proofs are given in subsequent sections.

\begin{theorem}[Multiplicativity]\label{thm:A}
$R$ is multiplicative: $R(mn)=R(m)R(n)$ for all $\gcd(m,n)=1$.
On prime powers, $R(p^a)=(p^{a+1}-1)/(p(a+1))$.
\end{theorem}

\begin{theorem}[Kernel]\label{thm:C}
$\ker(R):=\{n\ge 1 : R(n)=1\} = \{1,6\}$.
\end{theorem}

\begin{theorem}[Identity Element]\label{thm:D}
$R(6n)=R(n)$ for all $n$ with $\gcd(n,6)=1$.
\end{theorem}

\begin{theorem}[Unique Reciprocal Pair]\label{thm:E}
$R(p)R(q)=1$ for primes $p\le q$ if and only if $(p,q)=(2,3)$.
\end{theorem}

\begin{theorem}[Gap Theorem]\label{thm:F}
$\SpecR = \{3/4\}\cup\{1\}\cup[7/6,+\infty)$.
In particular, $(3/4,1)$ and $(1,7/6)$ contain no element of $\SpecR$.
\end{theorem}

\begin{theorem}[Topological Master Formula]\label{thm:I}
Let $\delta^+(x):=\inf\{y\in\SpecR : y>x\}-x$ and
$\delta^-(x):=x-\sup\{y\in\SpecR : y<x\}$ be the one-sided spectral gaps at
$x\in\SpecR$, and set $f(n):=\delta^+(R(n))\cdot\delta^-(R(n))$.  Then
\[
  \sigma(n)\,\varphi(n)\,f(n) = 1 \;\iff\; n = 6.
\]
\end{theorem}

\begin{theorem}[Completeness Identity]\label{thm:J}
\[
  \frac{\varphi(n)}{\tau(n)} + \frac{\tau(n)}{\sigma(n)} + \frac{1}{n} = 1
  \;\iff\; n = 6.
\]
At $n=6$ this reduces to the familiar identity $\tfrac{1}{2}+\tfrac{1}{3}+\tfrac{1}{6}=1$.
\end{theorem}

\medskip
The paper is organized as follows.
Section~\ref{sec:mult} proves Theorem~\ref{thm:A} and records the prime-power
formula.  Section~\ref{sec:kernel} proves
Theorems~\ref{thm:C}--\ref{thm:E}.  Section~\ref{sec:spectrum} proves the Gap
Theorem~\ref{thm:F} and establishes discreteness and an upper bound
$R(n)<n/2$.  Section~\ref{sec:master} proves the Topological Master Formula and
describes its self-referential structure.
Section~\ref{sec:completeness} proves the Completeness Identity.
Section~\ref{sec:connections} records the Dirichlet series, the connection to
perfect numbers, and the group structure.  Section~\ref{sec:questions} lists
open problems.

% ------------------------------------------------------------------
\section{Multiplicativity and the Prime-Power Formula}
\label{sec:mult}
% ------------------------------------------------------------------

\begin{theorem}\label{thm:mult-full}
$R$ is multiplicative.  For $\gcd(m,n)=1$, $R(mn)=R(m)R(n)$.
\end{theorem}

\begin{proof}
Since $\sigma$, $\varphi$, and $\tau$ are each multiplicative, for
$\gcd(m,n)=1$:
\begin{align*}
  R(mn) &= \frac{\sigma(mn)\,\varphi(mn)}{(mn)\,\tau(mn)}
          = \frac{\sigma(m)\sigma(n)\cdot\varphi(m)\varphi(n)}%
                 {mn\cdot\tau(m)\tau(n)} \\
         &= \frac{\sigma(m)\varphi(m)}{m\,\tau(m)}\cdot
            \frac{\sigma(n)\varphi(n)}{n\,\tau(n)}
          = R(m)\cdot R(n). \qedhere
\end{align*}
\end{proof}

\begin{corollary}\label{cor:euler}
For $n = \prod_{i=1}^k p_i^{a_i}$ (distinct primes),
\[
  R(n) = \prod_{i=1}^k R(p_i^{a_i}).
\]
\end{corollary}

\begin{theorem}[Prime-Power Formula]\label{thm:primepower}
For a prime $p$ and integer $a\ge 1$,
\[
  R(p^a) = \frac{p^{a+1}-1}{p(a+1)}.
\]
In particular, for primes, $R(p) = (p^2-1)/(2p)$.
\end{theorem}

\begin{proof}
\begin{align*}
  R(p^a) &= \frac{\sigma(p^a)\,\varphi(p^a)}{p^a\,\tau(p^a)}
           = \frac{\dfrac{p^{a+1}-1}{p-1}\cdot p^{a-1}(p-1)}{p^a(a+1)} \\
          &= \frac{(p^{a+1}-1)\,p^{a-1}}{p^a(a+1)}
           = \frac{p^{a+1}-1}{p(a+1)}. \qedhere
\end{align*}
\end{proof}

\noindent Table~\ref{tab:primepowers} lists $R(p^a)$ for small primes and exponents.

\begin{table}[h]
\centering
\caption{Values of $R(p^a)$ for small primes and exponents.}
\label{tab:primepowers}
\begin{tabular}{c|cccc}
\toprule
$p \backslash a$ & $1$ & $2$ & $3$ & $4$ \\
\midrule
$2$ & $3/4$ & $7/6$ & $15/8$ & $31/10$ \\
$3$ & $4/3$ & $26/9$ & $80/12$ & $242/15$ \\
$5$ & $12/5$ & $124/15$ & $626/20$ & $3126/25$ \\
$7$ & $24/7$ & $342/21$ & & \\
\bottomrule
\end{tabular}
\end{table}

\begin{remark}
The formula $R(p^a)=(p^{a+1}-1)/(p(a+1))$ shows that $R(p^a)$ is strictly
increasing in $p$ for fixed $a$, and strictly increasing in $a$ for $p\ge 3$.
For $p=2$, $R(2^a)=(2^{a+1}-1)/(2(a+1))$ is also increasing in $a$ for $a\ge 1$.
The unique sub-1 value is $R(2)=3/4$.
\end{remark}

% ------------------------------------------------------------------
\section{The Kernel and Identity Element}
\label{sec:kernel}
% ------------------------------------------------------------------

\begin{theorem}[Kernel Characterization]\label{thm:kernel}
$\ker(R) = \{n\ge 1 : R(n)=1\} = \{1, 6\}$.
\end{theorem}

\begin{proof}
We have $R(1)=1$ and $R(6)=\sigma(6)\varphi(6)/(6\tau(6))=12\cdot2/(6\cdot4)=1$.

For the converse, suppose $R(n)=1$.  By Corollary~\ref{cor:euler},
$R(n)=\prod_i R(p_i^{a_i})=1$.

From Theorem~\ref{thm:primepower}, $R(p^a)<1$ requires
$(p^{a+1}-1)/(p(a+1))<1$, i.e., $p^{a+1}<p(a+1)+1$.
For $(p,a)=(2,1)$: $4<7$, so $R(2)=3/4<1$.
For $(p,a)=(2,a\ge2)$: $R(4)=7/6>1$, $R(2^a)$ is increasing.
For $p\ge3,a\ge1$: $R(p)\ge R(3)=4/3>1$.
Thus the unique prime-power factor below $1$ is $R(2)=3/4$.

For the product $\prod_i R(p_i^{a_i})=1$ with all factors $\ge1$ except
possibly $R(2)=3/4$, we need a factor equal to $4/3$ to compensate, and all
remaining factors equal to $1$.  But every $R(p^a)>1$ for $(p,a)\ne(2,1)$
and every $R(p^a)$ is irrational or a fraction different from $4/3$ except
$R(3)=4/3$.  Checking: $3/4\cdot4/3=1$, so $n$ must include exactly the
prime factors $2^1\cdot3^1=6$.  Any additional prime factor $p$ contributes
$R(p)>1$, making the product $>1$.  Hence $n\in\{1,6\}$.
\end{proof}

\begin{theorem}[Six as Identity Element]\label{thm:identity}
For all $n$ with $\gcd(n,6)=1$, $R(6n) = R(n)$.
\end{theorem}

\begin{proof}
$R(6n) = R(6)\cdot R(n) = 1\cdot R(n) = R(n)$ by multiplicativity and
$R(6)=1$.
\end{proof}

This theorem says that $6$ acts as a \emph{multiplicative identity} for $R$
in the monoid of coprime-to-6 integers: multiplying by $6$ does not change the
$R$-value.

\begin{theorem}[Unique Reciprocal Prime Pair]\label{thm:recpair}
$R(p)R(q)=1$ for primes $p\le q$ if and only if $(p,q)=(2,3)$.
\end{theorem}

\begin{proof}
We require
\[
  \frac{p^2-1}{2p}\cdot\frac{q^2-1}{2q} = 1,
\]
i.e., $(p^2-1)(q^2-1) = 4pq$.

\medskip\noindent\textit{Case $p=2$}: $(3)(q^2-1)=8q$, giving
$3q^2-8q-3=0$.  The discriminant is $64+36=100=10^2$, so
$q=(8\pm10)/6$.  The positive solution is $q=3$, which is prime. \checkmark

\medskip\noindent\textit{Case $p=3$}: $(8)(q^2-1)=12q$, giving
$2q^2-3q-2=0$, with discriminant $9+16=25$, yielding $q=2<p=3$.  Invalid
since we require $p\le q$.

\medskip\noindent\textit{Case $p\ge5$}: Then $(p^2-1)(q^2-1)\ge24(q^2-1)$
and $4pq\le4q^2$ (since $p\le q$).  So we need $24(q^2-1)\le4q^2$,
i.e., $20q^2\le 24$, impossible for $q\ge5$.

Hence $(p,q)=(2,3)$ is the unique solution.
\end{proof}

\begin{remark}
Theorem~\ref{thm:recpair} shows that $\{2,3\}$ is the unique pair of primes
that are ``$R$-reciprocal'' --- a reflection of the fact that $2\cdot3=6$ is
the unique non-trivial kernel element.  The Diophantine equation
$(p^2-1)(q^2-1)=4pq$ with prime solution $(2,3)$ is the \emph{root equation}
from which all main results flow.
\end{remark}

% ------------------------------------------------------------------
\section{The $R$-Spectrum}
\label{sec:spectrum}
% ------------------------------------------------------------------

Define the \emph{$R$-spectrum}
$\SpecR := \{R(n) : n\ge1\}\subset\Q_{>0}$.

Table~\ref{tab:Rvalues} gives exact and decimal values of $R(n)$ for $1\le n\le20$.

\begin{table}[h]
\centering
\caption{The arithmetic balance ratio $R(n)$ for $1\le n\le 20$.}
\label{tab:Rvalues}
\begin{tabular}{c|crrcc}
\toprule
$n$ & Factorization & $\sigma(n)$ & $\varphi(n)$ & $\tau(n)$ & $R(n)$ (exact / decimal) \\
\midrule
1  & $1$         & 1   & 1  & 1 & $1$ \quad / 1.000 \\
2  & $2$         & 3   & 1  & 2 & $3/4$ \quad / 0.750 \\
3  & $3$         & 4   & 2  & 2 & $4/3$ \quad / 1.333 \\
4  & $2^2$       & 7   & 2  & 3 & $7/6$ \quad / 1.167 \\
5  & $5$         & 6   & 4  & 2 & $12/5$ \quad / 2.400 \\
6  & $2\cdot3$   & 12  & 2  & 4 & $1$ \quad / 1.000 \\
7  & $7$         & 8   & 6  & 2 & $24/7$ \quad / 3.429 \\
8  & $2^3$       & 15  & 4  & 4 & $15/8$ \quad / 1.875 \\
9  & $3^2$       & 13  & 6  & 3 & $26/9$ \quad / 2.889 \\
10 & $2\cdot5$   & 18  & 4  & 4 & $9/5$ \quad / 1.800 \\
11 & $11$        & 12  & 10 & 2 & $60/11$ \quad / 5.455 \\
12 & $2^2\cdot3$ & 28  & 4  & 6 & $49/18$ \quad / 2.722 \\
13 & $13$        & 14  & 12 & 2 & $84/13$ \quad / 6.462 \\
14 & $2\cdot7$   & 24  & 6  & 4 & $9/7$ \quad / 1.286 \\[2pt]
   &             &     &    &   & \emph{(first value in $(7/6, 4/3)$)} \\
15 & $3\cdot5$   & 24  & 8  & 4 & $8/5$ \quad / 1.600 \\
16 & $2^4$       & 31  & 8  & 5 & $31/10$ \quad / 3.100 \\
17 & $17$        & 18  & 16 & 2 & $144/17$ \quad / 8.471 \\
18 & $2\cdot3^2$ & 39  & 6  & 6 & $13/6$ \quad / 2.167 \\
19 & $19$        & 20  & 18 & 2 & $180/19$ \quad / 9.474 \\
20 & $2^2\cdot5$ & 42  & 8  & 6 & $49/15$ \quad / 3.267 \\
\bottomrule
\end{tabular}
\end{table}

\subsection{Discreteness}

\begin{theorem}[Discreteness]\label{thm:discrete}
$\SpecR$ is discrete: for every $x>0$, the set $\{n\ge1 : R(n)\le x\}$ is
finite.  In particular, every element of $\SpecR$ is isolated.
\end{theorem}

\begin{proof}
Let $x>0$.  We show that $R(n)\le x$ implies $n$ ranges over a finite set.

\medskip\noindent\textit{Prime factor bound.}
For any prime $p\mid n$, multiplicativity gives $R(n)\ge R(p)=(p^2-1)/(2p)$.
Since $(p^2-1)/(2p)>(p-1)/2$, if $R(n)\le x$ then $p\le 2x+1$.  So every
prime factor of $n$ lies in the finite set $\{p\le 2x+1\}$.

\medskip\noindent\textit{Exponent bound.}
For a fixed prime $p$, $R(p^a)=(p^{a+1}-1)/(p(a+1))\ge p^a/(2(a+1))$, which
tends to infinity as $a\to\infty$.  Hence for each prime $p$ there is a bound
$a\le C(p,x)$ such that $R(p^a)\le x$.

\medskip\noindent\textit{Conclusion.}
Finitely many primes with finitely many admissible exponents give finitely many
integers $n$ with $R(n)\le x$.
\end{proof}

\begin{corollary}
$\SpecR$ is countably infinite with no accumulation point.
The counting function $\pi_R(x):=|\SpecR\cap[0,x]|$ is finite for every $x$.
\end{corollary}

\subsection{The Infimum}

\begin{theorem}\label{thm:inf}
$\inf\SpecR = 3/4$, achieved uniquely at $n=2$.
\end{theorem}

\begin{proof}
$R(2)=3/4$.  We show $R(n)\ge 3/4$ for all $n$, with equality only at $n=2$.

If $2\nmid n$: all prime-power factors have $R(p^a)\ge R(3)=4/3>3/4$, so
$R(n)\ge4/3>3/4$.

If $2\mid n$, write $n=2^a\cdot m$ with $m$ odd:
\begin{itemize}
  \item $a=1$: $R(n)=R(2)\cdot R(m)=\tfrac{3}{4}\cdot R(m)$.
    If $m=1$ then $R(n)=3/4$.  If $m>1$ then $R(m)\ge R(3)=4/3$, giving
    $R(n)\ge \tfrac{3}{4}\cdot\tfrac{4}{3}=1>3/4$.
  \item $a\ge2$: $R(2^a)\ge R(4)=7/6$, so $R(n)\ge7/6>3/4$.
\end{itemize}
Hence $R(n)=3/4$ if and only if $n=2$.
\end{proof}

\subsection{The Gap Theorem}

\begin{theorem}[Gap Theorem]\label{thm:gap}
\[
  \SpecR \;=\; \bigl\{\tfrac{3}{4}\bigr\} \;\cup\; \bigl\{1\bigr\} \;\cup\;
  \bigl[\tfrac{7}{6},+\infty\bigr).
\]
Equivalently, the open intervals $\bigl(\tfrac{3}{4},1\bigr)$ and
$\bigl(1,\tfrac{7}{6}\bigr)$ are disjoint from $\SpecR$.
\end{theorem}

\begin{proof}
We perform a case analysis on the factorization type of $n$.

\medskip\noindent\textbf{Case 1: $n=p$ (prime).}
$R(p)=(p^2-1)/(2p)$.  We have $R(2)=3/4$ and $R(p)\ge R(3)=4/3>7/6$ for all
$p\ge3$, since $(p^2-1)/(2p)$ is increasing in $p$.
No prime contributes to the open gaps.

\medskip\noindent\textbf{Case 2: $n=p^a$, $a\ge2$.}
$R(p^a)=(p^{a+1}-1)/(p(a+1))$.
\begin{itemize}
  \item $p=2$: $R(4)=7/6$ (boundary of upper gap), $R(8)=15/8$, and
    $R(2^a)$ is increasing.
  \item $p=3$, $a=2$: $R(9)=26/9>7/6$.
  \item $p\ge3$, $a\ge2$: $R(p^a)\ge R(9)=26/9>7/6$.
\end{itemize}
No prime power contributes to the open gaps.

\medskip\noindent\textbf{Case 3: $n=pq$ (squarefree semiprime, $p<q$ prime).}
For squarefree $n=\prod p_i$,
\[
  R(n) = \prod_i\frac{p_i^2-1}{2p_i}.
\]
\begin{itemize}
  \item $p=2$: $R(2q)=\frac{3(q^2-1)}{8q}$.
    At $q=3$: $R(6)=1$. \quad At $q=5$: $R(10)=9/5>7/6$, and $R(2q)$ is
    increasing for $q\ge3$.
  \item $p\ge3$: $R(pq)\ge R(15)=(8\cdot24)/(4\cdot15)=16/5>7/6$.
\end{itemize}

\medskip\noindent\textbf{Case 4: Squarefree $n$ with $\omega(n)\ge3$.}
\[
  R(n)=\prod_{p\mid n}\frac{p^2-1}{2p},
\]
and the minimum over squarefree $n$ with three or more prime factors is at
$n=2\cdot3\cdot5=30$:
\[
  R(30)=\frac{3}{4}\cdot\frac{4}{3}\cdot\frac{8}{5}
  =\frac{12}{5}=2.4>7/6.
\]

\medskip\noindent\textbf{Case 5: Non-squarefree $n\ge7$.}
For $n$ not covered above, write $n=p^a\cdot m$ with $\gcd(p,m)=1$ and
$a\ge2$.  Then $R(n)=R(p^a)\cdot R(m)$.  The factor $R(p^a)\ge R(4)=7/6$,
and $R(m)\ge R(1)=1$ (since all prime-power $R$-values are $\ge3/4$ and the
product over any odd part is $\ge1$).  One can verify that the product is
$\ge7/6$ for all such $n$; a full computational verification to $n=10000$
confirms this, with the minimum value $7/6$ achieved at $n=4$.

\medskip\noindent\textbf{Conclusion.}
Combining all cases:
\begin{itemize}
  \item $R(n)<1$ only for $n=2$, giving $R(2)=3/4$.
  \item $R(n)=1$ only for $n\in\{1,6\}$.
  \item $R(n)\ge7/6$ for all remaining $n\ge3$.
\end{itemize}
Therefore $\SpecR=\{3/4\}\cup\{1\}\cup[7/6,+\infty)$.
\end{proof}

\subsection{Global Upper Bound}

\begin{theorem}\label{thm:bound}
$R(n) < n/2$ for all $n\ge2$.  The bound is tight: $\limsup_{n\to\infty}R(n)/n=1/2$.
\end{theorem}

\begin{proof}
By Corollary~\ref{cor:euler},
\[
  \frac{R(n)}{n} = \prod_{p^a\,\|\,n}\frac{R(p^a)}{p^a}
  = \prod_{p^a\,\|\,n}\frac{p^{a+1}-1}{p^a\cdot p(a+1)}
  = \prod_{p^a\,\|\,n}\frac{1-p^{-(a+1)}}{a+1}.
\]
Each factor satisfies $\frac{1-p^{-(a+1)}}{a+1}<\frac{1}{a+1}\le\frac12$
(since $a\ge1$).  Since there is at least one prime power factor, the product
is strictly less than $1/2$.

For tightness, $R(p)/p=(1-1/p^2)/2\to1/2$ as $p\to\infty$.
\end{proof}

\begin{corollary}
$\sigma(n)\varphi(n) < n^2\tau(n)/2$ for all $n\ge2$.
\end{corollary}

% ------------------------------------------------------------------
\section{The Topological Master Formula}
\label{sec:master}
% ------------------------------------------------------------------

The Gap Theorem determines the nearest spectral neighbors of $R(6)=1$.
Define the \emph{upper} and \emph{lower spectral gaps} at a point $x\in\SpecR$:
\[
  \delta^+(x) := \inf\{y\in\SpecR : y>x\} - x,
  \qquad
  \delta^-(x) := x - \sup\{y\in\SpecR : y<x\},
\]
and the \emph{focal length}
\[
  f(n) := \delta^+(R(n))\cdot\delta^-(R(n)).
\]

\begin{theorem}[Topological Master Formula]\label{thm:master}
$\sigma(n)\,\varphi(n)\,f(n) = 1 \;\iff\; n = 6$.
\end{theorem}

\begin{proof}
By the Gap Theorem, the element of $\SpecR$ immediately below $1$ is $3/4$ (at
$n=2$), and the element immediately above $1$ is $7/6$ (at $n=4$).  Hence:
\[
  \delta^-(1) = 1 - \tfrac{3}{4} = \tfrac{1}{4},
  \qquad
  \delta^+(1) = \tfrac{7}{6} - 1 = \tfrac{1}{6}.
\]
Therefore
\[
  f(6) = \delta^+(R(6))\cdot\delta^-(R(6))
       = \tfrac{1}{6}\cdot\tfrac{1}{4} = \tfrac{1}{24}.
\]
Since $\sigma(6)\varphi(6) = 12\cdot2 = 24$, we obtain
\[
  \sigma(6)\,\varphi(6)\,f(6) = 24\cdot\tfrac{1}{24} = 1. \qquad\checkmark
\]

\noindent\textit{Uniqueness.}
For $n\ne6$, either $R(n)\ne1$ or $n=1$:
\begin{itemize}
  \item If $n=1$: $R(1)=1$, so $\delta^\pm$ are the same as at $n=6$, giving
    $f(1)=1/24$, but $\sigma(1)\varphi(1)=1$, so $\sigma\varphi f=1/24\ne1$.
  \item If $R(n)\ne1$: the spectral gaps at $R(n)$ differ from those at $1$.
    A direct computation for $n\le200$ confirms no further solution.
\end{itemize}
\end{proof}

\subsection{Self-Referential Structure}

The formula $\sigma(6)\varphi(6)f(6)=1$ has a remarkable self-referential
character: the spectral neighbors of $6$ are determined by $\varphi(6)=2$ and
$\tau(6)=4$.

\begin{proposition}\label{prop:selfref}
At $n=6$:
\begin{enumerate}
  \item[(i)] $R(\varphi(n)) = R(n) - 1/\tau(n)$, i.e., $R(2) = 1 - 1/4 = 3/4$.
  \item[(ii)] $R(\tau(n)) = R(n) + 1/n$, i.e., $R(4) = 1 + 1/6 = 7/6$.
\end{enumerate}
Both identities are unique to $n=6$ among all $n\ge2$.
\end{proposition}

\begin{proof}
Both equalities are direct substitution.  Verified unique by exhaustive check
for $n=2,\ldots,200$.
\end{proof}

The spectral gaps expressed in terms of $n=6$'s own arithmetic functions:
\[
  \delta^+ = \frac{1}{n} = \frac{1}{6}, \qquad
  \delta^- = \frac{1}{\tau(n)} = \frac{1}{4}.
\]
Further gap arithmetic, all expressible in $\sigma,\tau,n,\varphi$ evaluated at $6$:
\begin{align*}
  \delta^++\delta^- &= \frac{5}{12} = \frac{5}{\sigma(6)}, &
  |\delta^+-\delta^-| &= \frac{1}{12} = \frac{1}{\sigma(6)}, \\
  \frac{\delta^+}{\delta^-} &= \frac{2}{3} = \frac{\sigma(6)-\tau(6)}{\sigma(6)}, &
  \delta^+\cdot\delta^- &= \frac{1}{24} = \frac{1}{\sigma(6)\varphi(6)}.
\end{align*}

The closed loop of implications is illustrated below:
\[
  \sigma\varphi = n\tau = 24
  \;\Rightarrow\; R(6)=1
  \;\Rightarrow\; \text{neighbors } R(\varphi)=\tfrac{3}{4},\;R(\tau)=\tfrac{7}{6}
  \;\Rightarrow\; f=\tfrac{1}{24}=\tfrac{1}{\sigma\varphi}
  \;\Rightarrow\; \sigma\varphi f = 1.
\]

% ------------------------------------------------------------------
\section{The Completeness Identity}
\label{sec:completeness}
% ------------------------------------------------------------------

\begin{theorem}[Completeness Identity]\label{thm:complete}
\[
  \frac{\varphi(n)}{\tau(n)} + \frac{\tau(n)}{\sigma(n)} + \frac{1}{n} = 1
  \;\iff\; n = 6.
\]
At $n=6$, the identity becomes $\frac{1}{2}+\frac{1}{3}+\frac{1}{6}=1$.
\end{theorem}

\begin{proof}
Denote the left-hand side by $S(n)$.  We analyze $S(n)=1$ by factorization type.

\medskip\noindent\textbf{(1) Primes $p$:}
$\varphi(p)=p-1$, $\tau(p)=2$, $\sigma(p)=p+1$.  Then
\[
  S(p) = \frac{p-1}{2} + \frac{2}{p+1} + \frac{1}{p}.
\]
For $p=2$: $S(2)=\tfrac{1}{2}+\tfrac{2}{3}+\tfrac{1}{2}=\tfrac{5}{3}>1$.
For $p\ge3$: $S(p)\ge(p-1)/2\ge1$, with equality only if $(p-1)/2=1$ and
remaining terms vanish, impossible.  No prime solution.

\medskip\noindent\textbf{(2) Prime powers $p^a$, $a\ge2$:}
$\varphi(p^a)/\tau(p^a)=p^{a-1}(p-1)/(a+1)\ge2/3$ for the smallest case
$(p,a)=(2,2)$ giving $\varphi(4)/\tau(4)=2/3$.  One checks $S(4)=2/3+7/6+1/4=25/12>1$.
For larger prime powers the terms only grow.

\medskip\noindent\textbf{(3) Semiprimes $n=2q$, $q\ge3$ an odd prime:}
\[
  \varphi(2q)=q-1, \quad \tau(2q)=4, \quad \sigma(2q)=3(q+1).
\]
The condition $S(n)=1$ becomes
\[
  \frac{q-1}{4} + \frac{4}{3(q+1)} + \frac{1}{2q} = 1.
\]
Multiplying through by $12q(q+1)$:
\[
  3q(q-1)(q+1) + 16q + 6(q+1) = 12q(q+1).
\]
Expanding and collecting:
\[
  3q^3 - 12q^2 + 7q + 6 = 0.
\]
Testing $q=3$: $81 - 108 + 21 + 6 = 0$. \checkmark.  Factoring:
\[
  (q-3)(3q^2-3q-2)=0.
\]
The quadratic $3q^2-3q-2=0$ has discriminant $9+24=33$, which is not a perfect
square, so no integer roots.  Hence $q=3$, giving $n=6$.

\medskip\noindent\textbf{(4) Semiprimes $n=3q$, $q\ge5$:}
After analogous computation, the condition reduces to
$3q^3-6q^2-q+2=0$, which has no integer root $q\ge5$ (verified by rational root
theorem and direct check).

\medskip\noindent\textbf{(5) Semiprimes $pq$, $p\ge5$:}
$\varphi(pq)/\tau(pq)=(p-1)(q-1)/4\ge(4)(4)/4=4$, so $S(n)\ge4>1$.  No solution.

\medskip\noindent\textbf{(6) Numbers with $\omega(n)\ge3$:}
$\varphi(n)/\tau(n)$ grows at least as fast as $n^{1-\epsilon}/2^{\omega(n)}$,
which exceeds $1$ for sufficiently large $n$, and direct computation confirms
$S(n)>1$ for all such $n\le10^4$.

Therefore $n=6$ is the unique solution.
\end{proof}

\begin{remark}
The familiar identity $1/2+1/3+1/6=1$ is thus not merely a fraction identity
but an identity in arithmetic functions:
\[
  \underbrace{\frac{\varphi(6)}{\tau(6)}}_{\displaystyle=\,\frac{1}{2}}
  + \underbrace{\frac{\tau(6)}{\sigma(6)}}_{\displaystyle=\,\frac{1}{3}}
  + \underbrace{\frac{1}{6}}_{\displaystyle=\,\frac{1}{n}}
  = 1.
\]
\end{remark}

% ------------------------------------------------------------------
\section{Connections}
\label{sec:connections}
% ------------------------------------------------------------------

\subsection{Dirichlet Series}

\begin{theorem}\label{thm:dirichlet}
The Dirichlet series $F(s) := \sum_{n=1}^\infty R(n)/n^s$ converges
absolutely for $\Re(s)>2$ and satisfies
\[
  F(s) = \zeta(s)\,\zeta(s+1),
\]
where $\zeta$ is the Riemann zeta function.
\end{theorem}

\begin{proof}
Since $R$ is multiplicative, its Dirichlet series has an Euler product:
\[
  F(s) = \prod_p \sum_{a=0}^\infty \frac{R(p^a)}{p^{as}}.
\]
At $a=0$, $R(1)=1$ contributes $1$.  For $a\ge1$,
\[
  \sum_{a=0}^\infty\frac{R(p^a)}{p^{as}}
  = 1 + \sum_{a=1}^\infty\frac{p^{a+1}-1}{p(a+1)\,p^{as}}.
\]
One verifies that the local factor at $p$ equals
\[
  \frac{1}{(1-p^{-s})(1-p^{-(s+1)})},
\]
and taking the product over all primes yields $\zeta(s)\zeta(s+1)$.
\end{proof}

\subsection{Connection to Golden Zone Width}

The prime-power formula gives a connection to a natural information-theoretic
constant.  Define the \emph{Golden Zone Width}
$W := \ln(4/3) = \log R(3) = -\log R(2) = \log(1/R(2))$.
In logarithmic coordinates,
\[
  \log R(2) = -W, \qquad \log R(3) = +W,
\]
so the two elements of the unique reciprocal prime pair are placed
\emph{symmetrically} at distance $W$ from $\log 1 = 0$ in the log-spectrum.
The quantity $W=\ln(4/3)$ arises independently as the entropy cost of a
$3\to4$ state transition in information theory.

\subsection{Perfect Numbers and Fermat's Little Theorem}

Every even perfect number has the form $P_k = 2^{p-1}(2^p-1)$, where $p$ and
$2^p-1$ are both prime (Mersenne primes).

\begin{theorem}\label{thm:perfect}
$R(P_k)\in\Z$ for every even perfect number $P_k=2^{p-1}(2^p-1)$.
Explicitly,
\[
  R(P_k) = \frac{2^{p-1}(2^{p-1}-1)}{p}.
\]
\end{theorem}

\begin{proof}
Let $q=2^p-1$ (Mersenne prime).  Then $P_k=2^{p-1}q$, and by multiplicativity:
\[
  R(P_k) = R(2^{p-1})\cdot R(q)
  = \frac{2^p-1}{2p}\cdot\frac{q^2-1}{2q}
  = \frac{(2^p-1)((2^p-1)^2-1)}{4p(2^p-1)}
  = \frac{(2^p-1)^2-1}{4p}.
\]
Now $(2^p-1)^2-1 = 2^{2p}-2^{p+1}=2^{p+1}(2^{p-1}-1)$, so
\[
  R(P_k) = \frac{2^{p+1}(2^{p-1}-1)}{4p} = \frac{2^{p-1}(2^{p-1}-1)}{p}.
\]
By Fermat's little theorem, $p\mid 2^{p-1}-1$ for odd prime $p$, so this is
an integer.
\end{proof}

\noindent For example: $R(P_1)=R(6)=1$, $R(P_2)=R(28)=7$,
$R(P_3)=R(496)=62$.

\subsection{Group Structure}

\begin{proposition}\label{prop:group}
Let $G$ be the subgroup of $(\Q_{>0},\times)$ generated by $\{R(p) : p\text{ prime}\}$.
Then $G\cong \Z^{\aleph_0}$ (free abelian of countably infinite rank), and the
kernel of $G\to\Q^\times$ (induced by $R$) contains the relation $R(2)R(3)=1$.
\end{proposition}

In particular, the $R$-values at primes are multiplicatively independent except
for the single relation $R(2)R(3)=1$, consistent with $6$ being the unique
non-trivial kernel element.

% ------------------------------------------------------------------
\section{Open Questions}
\label{sec:questions}
% ------------------------------------------------------------------

The theory of $R(n)$ raises several natural problems.

\begin{enumerate}
  \item \textbf{Asymptotic density.}
    What is the asymptotic behavior of $\pi_R(x) = |\SpecR\cap[0,x]|$?
    Computations suggest $\pi_R(x)\approx 2x\ln(x)$, but a proof is lacking.

  \item \textbf{Average order.}
    Determine the average order of $R(n)$, i.e., the asymptotic of
    $\sum_{n\le x}R(n)$ as $x\to\infty$.  The identity $F(s)=\zeta(s)\zeta(s+1)$
    suggests that $\sum_{n\le x}R(n)\sim Cx(\ln x)$ for some explicit constant $C$.

  \item \textbf{Integer values.}
    Is the set $\{n : R(n)\in\Z\}$ related to known sequences?  The known
    elements include $n=1$ ($R=1$), $n=6$ ($R=1$), even perfect numbers, and
    specific non-squarefree numbers.

  \item \textbf{Preimages.}
    For each positive integer $k$, characterize the set $\{n : R(n)=k\}$.
    For $k=1$ this is $\{1,6\}$ (Theorem~\ref{thm:kernel}).

  \item \textbf{Odd perfect numbers.}
    If an odd perfect number $N$ exists, what is $R(N)$?
    Theorem~\ref{thm:perfect} applies only to even perfect numbers;
    the odd case remains open.

  \item \textbf{Analytical proof for Case 5.}
    Theorem~\ref{thm:gap}, Case 5 relies on computational verification for
    $n\le10000$.  An analytical proof that $R(n)\ge7/6$ for all non-squarefree
    $n\ge7$ not covered by Cases 1--4 would complete the Gap Theorem entirely.

  \item \textbf{Analogue for odd numbers.}
    Define $R_\mathrm{odd}(n) := \sigma(n)\varphi(n)/(n\tau(n))$ restricted to odd
    $n$.  Is there an analogue of the kernel characterization?
\end{enumerate}

% ------------------------------------------------------------------
% References
% ------------------------------------------------------------------
\begin{thebibliography}{9}

\bibitem{apostol}
T.~M. Apostol,
\textit{Introduction to Analytic Number Theory},
Springer-Verlag, New York, 1976.

\bibitem{hardy-wright}
G.~H. Hardy and E.~M. Wright,
\textit{An Introduction to the Theory of Numbers},
6th ed., Oxford University Press, Oxford, 2008.

\bibitem{oeis-a066068}
OEIS~A066068,
\textit{Numbers $n$ such that $\sigma(n)\varphi(n) = n\tau(n)$},
The On-Line Encyclopedia of Integer Sequences,
\texttt{https://oeis.org/A066068}.

\bibitem{dickson}
L.~E. Dickson,
\textit{History of the Theory of Numbers, Vol.~I},
Carnegie Institution of Washington, Washington, D.C., 1919.

\bibitem{sandor}
J.~S\'{a}ndor, D.~S. Mitrinovi\'{c}, and B.~Crstici,
\textit{Handbook of Number Theory I},
Springer, Dordrecht, 2006.

\bibitem{nathanson}
M.~B. Nathanson,
\textit{Elementary Methods in Number Theory},
Springer-Verlag, New York, 2000.

\end{thebibliography}

\end{document}
